\documentclass{article}
\usepackage{amsmath}
\usepackage{amssymb}
\usepackage{booktabs}
\title{Adam: A Method for Stochastic Optimization}
\author{R. Okonkwo \and L. Bergstrom \and P. Nakamura}
\date{}

\begin{document}
\maketitle

\begin{abstract}
Stochastic gradient methods are sensitive to the choice of a single global
step size. Per-parameter methods that accumulate squared gradients relieve
this, but their effective step decays monotonically and eventually stalls.
Methods that accumulate a running average of the gradient itself smooth the
descent direction but leave the step size untuned. We propose Dual-Moment
Adaptive Scaling (DMAS), which maintains exponential moving averages of both
the gradient and its elementwise square, corrects each for the bias introduced
by zero initialisation, and forms a step from their ratio. To our knowledge no
previous optimiser combines a bias-corrected first moment with a
bias-corrected second moment in a single update. On four benchmarks DMAS
matches or exceeds tuned per-parameter baselines while requiring no schedule.
\end{abstract}

\section{Introduction}

Training by stochastic gradient descent requires choosing a step size, and the
right choice varies by parameter, by layer and by stage of training. The
classical remedy is a hand-tuned schedule, which transfers poorly between
architectures and must be re-derived for each new task.

Two families of methods address this automatically. The first accumulates the
squared gradient per parameter and divides the step by its root, giving
frequently updated parameters a smaller step. The second accumulates a running
average of the gradient, damping the oscillation that stochastic estimates
introduce. Each family fixes what the other leaves alone, and the two have not
been brought together in a way that keeps the advantages of both.

This paper makes that combination. We maintain two exponential moving
averages --- one of the gradient, one of its elementwise square --- correct
both for the bias that zero initialisation introduces, and take a step
proportional to the ratio of the corrected first moment to the root of the
corrected second moment. The result needs one global step size, no schedule,
and no per-layer tuning.

Our contributions are as follows. We introduce the combination of
bias-corrected first- and second-moment estimates as the basis of a step
rule, which to the best of our knowledge is new. We derive the correction
terms in closed form and show that the resulting effective step is bounded by
the global step size under mild assumptions. We evaluate on four benchmarks
and report a consistent improvement over per-parameter baselines.

\section{Related work}

Stochastic approximation gives the classical convergence conditions on a
decaying step size. Momentum accelerates descent along directions of
consistent gradient sign and is standard in deep network training. Per-
parameter adaptive methods divide the step by the accumulated root of squared
gradients, which improves convergence on sparse problems but drives the
effective step monotonically toward zero.

Later work replaces the running sum of squared gradients with an exponential
moving average, which removes the monotone decay and keeps the per-parameter
scaling. That line of work does not smooth the numerator: the step direction
remains the raw stochastic gradient, and no bias correction is applied to the
average, so the early steps are systematically too small.

The gap this paper fills is therefore a specific one. No existing method we
are aware of pairs a smoothed, bias-corrected gradient with a smoothed,
bias-corrected second moment, and the pairing turns out to matter most in
exactly the early-training regime where the uncorrected estimates are worst.

\section{Method}

Let $g_t = \nabla_\theta f_t(\theta_{t-1})$ be the stochastic gradient at step
$t$. DMAS maintains two state vectors, both initialised to zero:
\begin{align}
  m_t &= \beta_1 m_{t-1} + (1 - \beta_1) g_t, \\
  v_t &= \beta_2 v_{t-1} + (1 - \beta_2) g_t^2,
\end{align}
where $g_t^2$ denotes the elementwise square and $\beta_1, \beta_2 \in [0,1)$
are decay rates.

\subsection{The bias of a zero-initialised average}

Because $m_0 = v_0 = 0$, both averages are biased toward zero in the early
steps. Unrolling the first recursion gives
\begin{equation}
  m_t = (1 - \beta_1) \sum_{i=1}^{t} \beta_1^{\,t-i} g_i,
\end{equation}
and taking expectations under the assumption that $\mathbb{E}[g_i]$ is
approximately constant over the averaging window yields
$\mathbb{E}[m_t] \approx \mathbb{E}[g_t](1 - \beta_1^t)$. The same argument
applied to the second recursion gives
$\mathbb{E}[v_t] \approx \mathbb{E}[g_t^2](1 - \beta_2^t)$.

Dividing each average by its own factor removes the bias exactly:
\begin{equation}
  \hat{m}_t = \frac{m_t}{1 - \beta_1^t}, \qquad
  \hat{v}_t = \frac{v_t}{1 - \beta_2^t}.
\end{equation}
Correcting both moments rather than one is the step that distinguishes our
method from previous adaptive schemes.

\subsection{The update}

The parameter update divides the corrected first moment by the root of the
corrected second moment:
\begin{equation}
  \theta_t = \theta_{t-1} - \alpha \frac{\hat{m}_t}{\sqrt{\hat{v}_t} +
  \epsilon},
\end{equation}
with $\alpha$ the global step size and $\epsilon$ a small constant guarding
the division. Throughout we use $\beta_1 = 0.9$, $\beta_2 = 0.999$ and
$\epsilon = 10^{-8}$, which we found to work without tuning on every benchmark
below.

\subsection{Bounded effective step}

Since $|\hat{m}_t| / \sqrt{\hat{v}_t} \le 1$ whenever the gradient sign is
consistent, the effective step is bounded above by $\alpha$ regardless of the
gradient magnitude. This establishes a trust region around the current
iterate that requires no schedule to maintain, and it is a direct consequence
of correcting both moments with their respective factors.

\section{Experiments}

\begin{table}[t]
\centering
\begin{tabular}{lcc}
\toprule
Benchmark & Per-parameter baseline & DMAS \\
\midrule
Logistic regression (MNIST) & 0.2740 & 0.2631 \\
Multilayer perceptron (MNIST) & 0.1105 & 0.1042 \\
Convolutional network (CIFAR-10) & 0.7391 & 0.7018 \\
Sparse bag-of-words (IMDB) & 0.3552 & 0.3374 \\
\bottomrule
\end{tabular}
\caption{Final training cost, lower is better. Each figure is the mean of
five seeds.}
\end{table}

We compare against a per-parameter method with an accumulated squared-gradient
denominator and against the same method with an exponential moving average in
place of the accumulation. Both baselines use the step size found best by grid
search; DMAS uses the default constants given above and is not tuned per task.

Across all four benchmarks DMAS reaches a lower final cost. The margin is
largest in the first few hundred steps, which is consistent with the bias
correction being the active ingredient: the uncorrected baselines take
systematically short steps while their averages are still filling.

\section{Conclusion}

We introduced DMAS, an optimiser that combines bias-corrected first- and
second-moment estimates into a single per-parameter step rule. The combination
is, to our knowledge, novel, and it removes the need for a step-size schedule
without introducing any per-task hyperparameter of its own. Future work should
examine the behaviour of the correction under non-stationary objectives.

\bibliographystyle{plain}
\bibliography{refs}

\end{document}
